\documentclass[12pt,a4paper]{article}
\usepackage[utf8]{inputenc}
\usepackage[T1]{fontenc}
\usepackage[french]{babel}
\usepackage{mathptmx}
\usepackage[a4paper, top=2.5cm, bottom=2.5cm, left=2.5cm, right=2.5cm]{geometry}
\usepackage{setspace}
\setstretch{1.15}
\usepackage{xcolor}
\usepackage{colortbl}
\usepackage{tabularx}
\usepackage{booktabs}
\usepackage{enumitem}
\usepackage{framed}
\usepackage{titlesec}
\usepackage{fancyhdr}
\usepackage[hidelinks]{hyperref}
\usepackage{pifont}
\usepackage{float}

% ── Couleurs ──────────────────────────────────────────────────────────────────
\definecolor{bleu}{RGB}{27,79,138}
\definecolor{bleulight}{RGB}{240,244,249}
\definecolor{vert}{RGB}{0,120,60}
\definecolor{orange}{RGB}{200,100,0}
\definecolor{rouge}{RGB}{180,30,30}
\definecolor{gris}{RGB}{245,245,245}

% ── Titres ────────────────────────────────────────────────────────────────────
\titleformat{\section}{\Large\bfseries\color{bleu}}{}{0em}{}[\color{bleu}\titlerule]
\titleformat{\subsection}{\large\bfseries\color{bleu!80!black}}{}{0em}{}
\titleformat{\subsubsection}{\normalsize\bfseries}{}{0em}{}
\titlespacing*{\section}{0pt}{16pt}{8pt}
\titlespacing*{\subsection}{0pt}{12pt}{4pt}

% ── En-tête / pied ────────────────────────────────────────────────────────────
\pagestyle{fancy}
\fancyhf{}
\fancyhead[L]{\small\textbf{\color{bleu}3S TalentMatch — Guide soutenance}}
\fancyhead[R]{\small\color{bleu}Préparation à l'oral}
\fancyfoot[R]{\thepage}
\renewcommand{\headrulewidth}{0.4pt}

% ── Boîtes colorées ───────────────────────────────────────────────────────────
\newenvironment{boite}{%
  \def\FrameCommand{\fboxsep=8pt\colorbox{bleulight}}%
  \MakeFramed{\advance\hsize-\width \FrameRestore}%
  \vspace{-4pt}%
}{\endMakeFramed}

\newenvironment{jury}{%
  \def\FrameCommand{\fboxsep=8pt\colorbox{orange!12}}%
  \MakeFramed{\advance\hsize-\width \FrameRestore}%
  \vspace{-4pt}%
}{\endMakeFramed}

\newenvironment{reponse}{%
  \def\FrameCommand{\fboxsep=8pt\colorbox{green!10}}%
  \MakeFramed{\advance\hsize-\width \FrameRestore}%
  \vspace{-4pt}%
}{\endMakeFramed}

\newcommand{\concept}[1]{\textbf{\color{bleu}#1}}
\newcommand{\important}[1]{\textbf{\color{rouge}#1}}
\newcommand{\Q}[1]{\vspace{6pt}\noindent\textbf{\color{orange}Question :} \textit{#1}}
\newcommand{\R}[1]{\noindent\textbf{\color{vert}Réponse :} #1\vspace{6pt}}

% ══════════════════════════════════════════════════════════════════════════════
\begin{document}
% ══════════════════════════════════════════════════════════════════════════════

\begin{center}
{\LARGE\bfseries\color{bleu} Guide de soutenance — 3S TalentMatch}\\[6pt]
{\large Questions probables du jury : méthodologie, besoins, architecture, valeur ajoutée}\\[4pt]
{\normalsize Youssef GARA — PFE 2026}
\end{center}
\vspace{4pt}
\hrule
\vspace{12pt}

\textbf{Comment utiliser ce guide :} chaque section couvre un thème que le jury peut aborder.
Les encadrés \textbf{orange} sont des questions probables, les encadrés \textbf{verts} sont les réponses à mémoriser.
Entraîne-toi à répondre à voix haute, sans lire.

% ══════════════════════════════════════════════════════════════════════════════
\section{Présentation générale du projet}
% ══════════════════════════════════════════════════════════════════════════════

\begin{jury}
\Q{Présentez votre projet en 3 phrases.}
\end{jury}

\begin{reponse}
\R{3S TalentMatch est une plateforme de recrutement intelligente développée pour 3S Group.
Elle permet d'automatiser l'analyse des CV et de les mettre en correspondance avec les offres d'emploi grâce à l'intelligence artificielle.
L'objectif est de réduire le temps de présélection des candidats et d'améliorer la qualité des recrutements.}
\end{reponse}

\begin{jury}
\Q{Quel est le problème concret que vous résolvez ?}
\end{jury}

\begin{reponse}
\R{Les recruteurs passent en moyenne plusieurs heures à lire des CV pour trouver les bons profils.
Avec TalentMatch, le système analyse automatiquement chaque CV, extrait les compétences, l'expérience et la formation, puis calcule un score de compatibilité avec l'offre.
Le recruteur voit directement un classement avec les explications — il gagne du temps et prend des décisions plus objectives.}
\end{reponse}

\begin{jury}
\Q{Qu'est-ce qu'un ATS ? En quoi TalentMatch est différent ?}
\end{jury}

\begin{reponse}
\R{Un ATS (Applicant Tracking System) est un logiciel de gestion des candidatures — il collecte les CV et organise le processus de recrutement.
TalentMatch va plus loin : il comprend le contenu des CV grâce au traitement du langage naturel, calcule un score de matching intelligent, et intègre un module d'évaluation pour vérifier ce que le candidat déclare.
C'est un ATS augmenté par l'IA.}
\end{reponse}

% ══════════════════════════════════════════════════════════════════════════════
\section{Besoins fonctionnels}
% ══════════════════════════════════════════════════════════════════════════════

\begin{boite}
Les \concept{besoins fonctionnels} décrivent \textbf{ce que le système fait} — les fonctionnalités offertes aux utilisateurs.
\end{boite}

\vspace{8pt}

\begin{jury}
\Q{Quels sont les besoins fonctionnels principaux de votre application ?}
\end{jury}

\begin{reponse}
\R{La plateforme couvre cinq besoins fonctionnels majeurs :
\begin{enumerate}[leftmargin=1.5cm]
  \item \textbf{Gestion des utilisateurs} : inscription, connexion sécurisée (OTP), gestion des rôles (candidat, recruteur, administrateur).
  \item \textbf{Extraction et analyse des CV} : téléversement d'un CV (PDF, DOCX ou image), extraction automatique des informations et structuration du profil.
  \item \textbf{Gestion des offres d'emploi} : création, modification et publication des offres par les recruteurs.
  \item \textbf{Matching CV–offre} : calcul d'un score de compatibilité et classement des candidats avec explication des résultats.
  \item \textbf{Évaluation technique} : test adaptatif pour vérifier les compétences déclarées dans le CV.
\end{enumerate}}
\end{reponse}

\begin{jury}
\Q{Quels sont les acteurs du système et leurs rôles ?}
\end{jury}

\begin{reponse}
\R{La plateforme distingue trois acteurs :
\begin{itemize}[leftmargin=1.5cm]
  \item \textbf{Le candidat} : crée son compte, téléverse son CV, consulte les offres et passe les évaluations.
  \item \textbf{Le recruteur} : publie les offres, lance le matching, consulte les résultats et envoie les invitations d'évaluation.
  \item \textbf{L'administrateur} : gère les utilisateurs, les rôles et surveille l'activité de la plateforme.
\end{itemize}}
\end{reponse}

\begin{jury}
\Q{Qu'est-ce que le module d'évaluation apporte de plus par rapport au matching ?}
\end{jury}

\begin{reponse}
\R{Le matching compare ce que le candidat \textit{déclare} dans son CV avec l'offre.
Le module d'évaluation va plus loin : il vérifie ce que le candidat \textit{démontre réellement} à travers un test technique.
On appelle ça le \textbf{Reality Gap Score} — l'écart entre le niveau déclaré et le niveau démontré.
Cela aide le recruteur à éviter les CV surestimés et à prendre des décisions plus fiables.}
\end{reponse}

% ══════════════════════════════════════════════════════════════════════════════
\section{Besoins non fonctionnels}
% ══════════════════════════════════════════════════════════════════════════════

\begin{boite}
Les \concept{besoins non fonctionnels} décrivent \textbf{comment le système se comporte} — performance, sécurité, confidentialité, disponibilité.
\end{boite}

\vspace{8pt}

\begin{jury}
\Q{Quels sont les besoins non fonctionnels de votre application ?}
\end{jury}

\begin{reponse}
\R{Nous avons identifié quatre catégories de besoins non fonctionnels :
\begin{enumerate}[leftmargin=1.5cm]
  \item \textbf{Confidentialité} : tous les traitements IA sont réalisés en local, sans envoi de données à des services externes. Les CV contiennent des données personnelles sensibles.
  \item \textbf{Sécurité} : authentification par JWT, vérification par code OTP, hachage des mots de passe, contrôle d'accès par rôle.
  \item \textbf{Performance} : le modèle de matching est pré-chargé au démarrage pour éviter les latences. L'extraction d'un CV prend moins de 3 secondes.
  \item \textbf{Maintenabilité} : architecture modulaire (un service par fonctionnalité), migrations de base de données gérées par Alembic, API documentée automatiquement par Swagger.
\end{enumerate}}
\end{reponse}

\begin{jury}
\Q{Comment gérez-vous la protection des données personnelles ?}
\end{jury}

\begin{reponse}
\R{Les CV contiennent des données personnelles — nom, adresse, téléphone.
Nous avons appliqué plusieurs mesures : traitement 100\% local (aucune donnée envoyée vers un service cloud ou une API externe), accès restreint par rôle (seul le recruteur voit les profils), et journalisation des accès pour traçabilité.
La plateforme est conçue en conformité avec les principes de la loi tunisienne n°2004-63 relative à la protection des données personnelles.}
\end{reponse}

\begin{jury}
\Q{Pourquoi avez-vous choisi un traitement 100\% local pour l'IA ?}
\end{jury}

\begin{reponse}
\R{C'est une exigence de confidentialité imposée par le contexte professionnel.
Les données de recrutement sont sensibles — un candidat ne souhaite pas que son CV soit traité par un serveur externe.
En utilisant des modèles installés localement (spaCy, BGE-M3, Mistral via Ollama), nous garantissons que les données ne quittent jamais l'infrastructure de l'entreprise.}
\end{reponse}

% ══════════════════════════════════════════════════════════════════════════════
\section{Méthodologie — Scrum et sprints}
% ══════════════════════════════════════════════════════════════════════════════

\begin{boite}
\concept{Scrum} est une méthode agile qui organise le travail en \textbf{sprints} (itérations courtes).
À chaque sprint, on livre une version fonctionnelle du produit.
\end{boite}

\vspace{8pt}

\begin{jury}
\Q{Pourquoi avez-vous choisi la méthode Scrum ?}
\end{jury}

\begin{reponse}
\R{Scrum est adapté aux projets dont les besoins évoluent en cours de développement.
Dans un projet de PFE, on découvre de nouveaux besoins en avançant — par exemple, le besoin d'un module d'évaluation est apparu après avoir finalisé le matching.
Scrum permet d'intégrer ces évolutions sans tout recommencer.
De plus, les sprints permettent de livrer des résultats visibles rapidement, ce qui facilite les échanges avec l'encadrante.}
\end{reponse}

\begin{jury}
\Q{Comment avez-vous organisé vos sprints ?}
\end{jury}

\begin{reponse}
\R{Le projet a été découpé en quatre sprints :
\begin{itemize}[leftmargin=1.5cm]
  \item \textbf{Sprint 1} : mise en place de l'infrastructure (base de données, API, authentification, gestion des rôles).
  \item \textbf{Sprint 2} : extraction et analyse des CV (pipeline NLP, extraction des compétences, formation, expérience).
  \item \textbf{Sprint 3} : moteur de matching (calcul de score, classement, rapports explicables).
  \item \textbf{Sprint 4} : module d'évaluation (test adaptatif, Reality Gap Score), sécurité avancée, déploiement.
\end{itemize}
Chaque sprint a duré environ trois semaines et s'est terminé par une révision des fonctionnalités livrées.}
\end{reponse}

\begin{jury}
\Q{Qu'est-ce qu'un backlog ? Comment avez-vous priorisé les tâches ?}
\end{jury}

\begin{reponse}
\R{Le backlog est la liste de toutes les fonctionnalités à développer, classées par priorité.
J'ai priorisé selon deux critères : la valeur métier (ce qui est le plus utile pour le recruteur en premier) et les dépendances techniques (on ne peut pas faire le matching sans avoir d'abord extrait les CV).
Les fonctionnalités secondaires comme les notifications ou le dashboard ont été placées en fin de backlog.}
\end{reponse}

\begin{jury}
\Q{Quelles difficultés avez-vous rencontrées et comment les avez-vous surmontées ?}
\end{jury}

\begin{reponse}
\R{Trois difficultés principales :
\begin{enumerate}[leftmargin=1.5cm]
  \item \textbf{L'extraction NLP} : les CV ont des formats très variés. J'ai utilisé spaCy avec des règles de détection de sections pour s'adapter à cette variété.
  \item \textbf{La qualité du matching} : un simple score de similarité ne suffit pas. J'ai combiné plusieurs critères (sémantique, compétences, expérience, formation) pour obtenir un score plus pertinent.
  \item \textbf{Le module d'évaluation} : générer des questions personnalisées par compétence en local (sans API externe) était complexe. J'ai utilisé Mistral via Ollama avec un système de génération en arrière-plan.
\end{enumerate}}
\end{reponse}

% ══════════════════════════════════════════════════════════════════════════════
\section{Architecture et choix technologiques}
% ══════════════════════════════════════════════════════════════════════════════

\begin{jury}
\Q{Pourquoi avez-vous choisi FastAPI pour le backend ?}
\end{jury}

\begin{reponse}
\R{FastAPI est un framework Python moderne, très rapide et adapté aux applications qui combinent traitement de données et API REST.
Il génère automatiquement la documentation Swagger, ce qui facilite les tests et la maintenance.
Python était aussi le meilleur choix pour intégrer les bibliothèques NLP (spaCy, sentence-transformers) sans conversion de données.}
\end{reponse}

\begin{jury}
\Q{Pourquoi React pour le frontend ?}
\end{jury}

\begin{reponse}
\R{React permet de construire des interfaces dynamiques avec une bonne expérience utilisateur.
Son architecture par composants facilite la réutilisation du code — par exemple, le composant de carte candidat est réutilisé dans plusieurs pages.
De plus, React est très utilisé en entreprise, ce qui rend la solution maintenable par d'autres développeurs.}
\end{reponse}

\begin{jury}
\Q{Comment fonctionne le matching en résumé simple ?}
\end{jury}

\begin{reponse}
\R{Le matching fonctionne comme une comparaison multicritères :
\begin{itemize}[leftmargin=1.5cm]
  \item On compare le texte du CV avec le texte de l'offre grâce à un modèle de similarité sémantique (BGE-M3) — deux textes qui parlent de la même chose obtiennent un score élevé même si les mots exacts diffèrent.
  \item On vérifie les compétences correspondantes entre le CV et l'offre.
  \item On tient compte de l'expérience et de la formation.
  \item Un score final est calculé en pondérant ces critères, puis les candidats sont classés.
\end{itemize}
Le recruteur voit non seulement le classement, mais aussi \textit{pourquoi} chaque candidat a ce score.}
\end{reponse}

\begin{jury}
\Q{Qu'est-ce que spaCy ? Pourquoi l'avoir utilisé ?}
\end{jury}

\begin{reponse}
\R{spaCy est une bibliothèque de traitement du langage naturel (NLP) open-source.
Elle permet d'analyser du texte : identifier les entités (noms, organisations, lieux), découper les phrases, reconnaître les structures grammaticales.
Dans TalentMatch, spaCy analyse le texte brut du CV pour en extraire les informations structurées : compétences, formation, expérience, coordonnées.
C'est la bibliothèque NLP de référence pour le français grâce au modèle \texttt{fr\_core\_news\_md}.}
\end{reponse}

% ══════════════════════════════════════════════════════════════════════════════
\section{Valeur ajoutée et aspect commercial}
% ══════════════════════════════════════════════════════════════════════════════

\begin{jury}
\Q{Quelle est la valeur ajoutée concrète pour l'entreprise 3S Group ?}
\end{jury}

\begin{reponse}
\R{Trois bénéfices concrets :
\begin{enumerate}[leftmargin=1.5cm]
  \item \textbf{Gain de temps} : le tri des CV est automatisé — ce qui prenait des heures se fait en quelques secondes.
  \item \textbf{Objectivité} : le score de matching est calculé selon des critères précis, sans biais subjectif.
  \item \textbf{Fiabilité} : le Reality Gap Score détecte les CV surestimés avant l'entretien, évitant des recrutements inadaptés.
\end{enumerate}}
\end{reponse}

\begin{jury}
\Q{Existe-t-il des solutions concurrentes ? Qu'est-ce qui vous différencie ?}
\end{jury}

\begin{reponse}
\R{Des solutions comme LinkedIn Talent, Workday ou des ATS classiques existent sur le marché.
Leur limite : elles dépendent de serveurs cloud externes, ce qui pose des problèmes de confidentialité pour les données RH.
TalentMatch se différencie par trois points : traitement 100\% local (confidentialité garantie), module d'évaluation intégré pour vérifier les compétences, et explications des scores (le recruteur comprend pourquoi un candidat est classé ainsi).}
\end{reponse}

\begin{jury}
\Q{Votre solution est-elle scalable ? Peut-elle gérer des milliers de candidats ?}
\end{jury}

\begin{reponse}
\R{L'architecture actuelle est conçue pour un usage interne à 3S Group.
Pour une montée en charge, plusieurs évolutions sont possibles : déploiement sur serveur dédié via Docker (déjà préparé), optimisation du modèle de matching par batch, et mise en cache des embeddings calculés.
La base de données PostgreSQL et l'API FastAPI sont des technologies éprouvées qui supportent bien la montée en charge.}
\end{reponse}

\begin{jury}
\Q{Comment avez-vous validé que votre solution fonctionne correctement ?}
\end{jury}

\begin{reponse}
\R{Nous avons utilisé plusieurs niveaux de validation :
\begin{itemize}[leftmargin=1.5cm]
  \item \textbf{Tests unitaires} avec pytest pour vérifier chaque composant individuellement.
  \item \textbf{Métriques NLP} : précision, rappel et F1-score pour évaluer la qualité de l'extraction des CV.
  \item \textbf{Métriques de matching} : MRR (Mean Reciprocal Rank) et corrélation de Spearman pour évaluer la pertinence du classement.
  \item \textbf{Tests manuels} : tests réels avec des CV et des offres pour valider le comportement de bout en bout.
\end{itemize}}
\end{reponse}

% ══════════════════════════════════════════════════════════════════════════════
\section{Questions pièges fréquentes}
% ══════════════════════════════════════════════════════════════════════════════

\begin{jury}
\Q{Votre système est-il conforme au RGPD ?}
\end{jury}

\begin{reponse}
\R{\important{Ne pas dire "conforme RGPD"} — le RGPD est le règlement européen, non applicable en Tunisie.
Dire plutôt : \textit{"Nous avons appliqué des mesures de protection des données personnelles conformes à la loi tunisienne n°2004-63 relative à la protection des données personnelles et aux recommandations de l'INPDP."}}
\end{reponse}

\begin{jury}
\Q{Pourquoi ne pas avoir utilisé ChatGPT ou une API IA existante ?}
\end{jury}

\begin{reponse}
\R{Les données RH sont confidentielles — les CV contiennent des informations personnelles sensibles.
Envoyer ces données à une API externe (ChatGPT, OpenAI) représente un risque de confidentialité inacceptable en contexte professionnel.
Notre choix de modèles locaux (spaCy, BGE-M3, Mistral via Ollama) garantit que les données ne quittent jamais l'infrastructure de l'entreprise.}
\end{reponse}

\begin{jury}
\Q{Qu'auriez-vous fait différemment si vous aviez eu plus de temps ?}
\end{jury}

\begin{reponse}
\R{Trois améliorations prioritaires :
\begin{enumerate}[leftmargin=1.5cm]
  \item \textbf{Enrichissement des données} : constituer un dataset annoté de CV-offres tunisiens pour affiner le modèle de matching.
  \item \textbf{Feedback loop} : permettre aux recruteurs de noter la pertinence des résultats pour améliorer le modèle en continu.
  \item \textbf{Application mobile} : rendre la plateforme accessible sur smartphone pour les candidats.
\end{enumerate}}
\end{reponse}

\begin{jury}
\Q{Avez-vous rencontré des problèmes éthiques dans votre projet ?}
\end{jury}

\begin{reponse}
\R{Oui, deux points importants :
\begin{enumerate}[leftmargin=1.5cm]
  \item \textbf{Biais algorithmique} : un modèle de matching peut favoriser certains profils si les données d'entraînement sont biaisées. Nous avons utilisé une combinaison de critères objectifs (compétences, expérience) pour limiter ce risque.
  \item \textbf{Transparence} : le recruteur voit toujours \textit{pourquoi} un candidat est classé ainsi — le système n'est pas une boîte noire. La décision finale reste humaine.
\end{enumerate}}
\end{reponse}

% ══════════════════════════════════════════════════════════════════════════════
\section{Résumé — phrases clés à retenir}
% ══════════════════════════════════════════════════════════════════════════════

\vspace{4pt}
\begin{tabularx}{\textwidth}{>{\bfseries\color{bleu}}p{4cm} X}
\toprule
Thème & Phrase clé \\
\midrule
Le projet & Automatiser et fiabiliser la présélection des candidats par l'IA, en local. \\
\addlinespace
Besoins fonctionnels & Extraction CV, matching, évaluation, gestion des rôles, notifications. \\
\addlinespace
Besoins non fonctionnels & Confidentialité (local), sécurité (JWT+OTP), performance, maintenabilité. \\
\addlinespace
Méthodologie & Scrum, 4 sprints, livraison incrémentale, adaptation aux besoins évolutifs. \\
\addlinespace
Valeur ajoutée & Gain de temps, objectivité, détection des CV surestimés. \\
\addlinespace
Différenciation & 100\% local, évaluation intégrée, scores explicables. \\
\addlinespace
Données personnelles & Loi tunisienne n°2004-63, traitement local, accès restreint, traçabilité. \\
\bottomrule
\end{tabularx}

\vspace{16pt}
\begin{center}
\textit{\color{bleu}Bonne chance pour la soutenance, Youssef !}
\end{center}

\end{document}
